\chapter{Le Fort I Osteotomy (Maxillary Osteotomy)}

\section*{What Is This Surgery?}
Le Fort I osteotomy is a critical surgical intervention aimed at correcting skeletal discrepancies of the maxilla, the upper jaw. This procedure involves the surgical separation of the maxilla from the craniofacial skeleton, allowing for its repositioning in three-dimensional space. The maxilla can be advanced, set back, impacted, or rotated to achieve optimal dental occlusion and improve facial aesthetics. This surgery is particularly significant in patients with malocclusion, facial asymmetry, or those requiring correction due to congenital conditions such as cleft lip and palate. Performed under general anesthesia by an oral and maxillofacial surgeon, Le Fort I osteotomy is often part of a comprehensive treatment plan that includes orthodontic therapy. The repositioned maxilla is stabilized using rigid internal fixation with titanium plates and screws, which facilitates healing and maintains the corrected position over time.

\section*{Alternative Names}
\begin{itemize}
  \item Maxillary osteotomy
  \item Upper jaw surgery
  \item Le Fort I advancement
  \item Maxillary impaction
  \item Orthognathic surgery of the maxilla
  \item Le Fort I procedure
\end{itemize}

\section*{Relevant Surgical Anatomy}
The Le Fort I osteotomy primarily involves the maxilla, a paired bone that forms the upper jaw and contributes to the structure of the face, nasal cavity, and orbit. Understanding the relevant anatomy is essential for minimizing complications during surgery. 

The osteotomy typically begins at the piriform aperture and extends laterally and posteriorly to the pterygomaxillary junction. Key anatomical structures include:

- **Infraorbital Nerve:** A branch of the trigeminal nerve (V2), it exits the infraorbital foramen and provides sensation to the midface. Care must be taken to avoid injury during dissection.
  
- **Greater Palatine Artery:** This artery supplies blood to the maxilla and runs along the hard palate. Its preservation is vital to maintain vascularity.
  
- **Maxillary Sinuses:** These air-filled cavities are located within the maxilla and are often entered during the osteotomy. Their integrity is crucial for preventing postoperative complications.
  
- **Pterygoid Plates:** Located posteriorly, these structures are important for the osteotomy and must be carefully navigated to avoid vascular injury.

- **Nasal Septum and Lateral Nasal Walls:** These structures are involved in mobilizing the maxilla and must be carefully managed during the procedure.

Understanding these relationships is critical for achieving successful surgical outcomes and minimizing complications.

\section*{Surgical Principle \& Rationale}
The primary surgical principle of the Le Fort I osteotomy is to correct dentofacial deformities resulting from an improperly positioned maxilla. Malocclusion and facial asymmetry often arise when the maxilla does not align properly with the mandible and cranial base. By performing a controlled horizontal osteotomy, the surgeon can mobilize the maxilla as a single unit or in segments, allowing for precise repositioning.

The rationale for this procedure is to restore optimal oral function, enhance facial aesthetics, and alleviate associated symptoms such as obstructive sleep apnea or speech difficulties. The technical goals include achieving a stable Class I occlusal relationship, improving facial aesthetics by correcting midface projection, and optimizing functional parameters such as mastication and nasal airflow. Rigid internal fixation with titanium plates and screws provides immediate stability, promoting primary bone healing and preventing relapse, thus ensuring long-term success of the surgical intervention.

\section*{History \& Surgical Evolution}
The Le Fort I osteotomy is named after the French surgeon René Le Fort, who, in 1901, described patterns of midfacial fractures based on cadaver studies. His work laid the groundwork for understanding facial anatomy and trauma, but it was not until the mid-20th century that these concepts were adapted for elective surgical procedures.

In the 1960s, Hugo Obwegeser refined the intraoral approach to maxillary repositioning, minimizing external scarring and improving surgical outcomes. William Bell's research in the 1970s on the vascular supply to the maxilla and the introduction of rigid internal fixation techniques revolutionized the procedure, enhancing safety and predictability. The advent of three-dimensional imaging and virtual surgical planning in recent years has further refined the Le Fort I osteotomy, allowing for precise preoperative planning and improved patient outcomes.

\section*{Clinical Indications}
The Le Fort I osteotomy is indicated for various skeletal and functional issues related to maxillary position, typically after skeletal growth has ceased. Indications include:

\begin{itemize}
  \item Correction of severe malocclusion, such as anterior open bite or deep overbite, when orthodontics alone is insufficient.
  \item Treatment of midface deficiency or hypoplasia, resulting in a concave facial profile.
  \item Improvement of facial aesthetics by correcting vertical maxillary excess or deficiency.
  \item Enhancement of chewing efficiency and masticatory function compromised by an unstable bite.
  \item Improvement of speech articulation affected by jaw position.
  \item Management of obstructive sleep apnea through maxillary advancement.
  \item Correction of facial asymmetry involving the maxilla.
  \item Addressing sequelae of cleft lip and palate, where maxillary growth is restricted.
\end{itemize}

\section*{Clinical Positioning}
The Le Fort I osteotomy is primarily a definitive surgical procedure for correcting skeletal maxillary discrepancies. It is typically performed after a period of pre-surgical orthodontics, which aligns the teeth and prepares the dental arches for surgical intervention. This initial orthodontic phase is crucial for ensuring that the teeth are positioned correctly over the basal bone, allowing for optimal surgical outcomes.

In certain cases, such as severe trauma or congenital deformities, the Le Fort I osteotomy may be performed as an immediate primary intervention. However, for elective corrections, it is usually integrated into a comprehensive treatment plan that includes pre-surgical orthodontics, the surgical phase, and post-surgical orthodontic finishing.

\section*{Surgical Technique \& Procedure}
The Le Fort I osteotomy is performed under general anesthesia, typically with nasotracheal intubation. The patient is positioned supine, and the surgical field is prepared in a sterile manner.

\begin{itemize}
  \item \textbf{Incision and Exposure:} An intraoral vestibular incision is made along the gum line, extending from the first molar region on one side to the other. A subperiosteal dissection exposes the maxilla, zygomatic buttresses, and piriform rims, while protecting the infraorbital nerves.
  
  \item \textbf{Osteotomy Cuts:} Using an oscillating saw, horizontal osteotomy cuts are made through the lateral maxillary walls, extending to the pterygomaxillary junction. Care is taken to avoid the apices of the maxillary teeth and the infraorbital foramen.
  
  \item \textbf{Downfracture and Mobilization:} The maxilla is downfractured and mobilized from the cranial base using specialized osteotomes, allowing for visualization of the nasal floor and posterior maxilla.
  
  \item \textbf{Repositioning and Occlusal Guidance:} The mobilized maxilla is repositioned into its planned position, guided by pre-fabricated surgical splints to ensure the desired occlusal relationship is achieved.
  
  \item \textbf{Rigid Internal Fixation:} The maxilla is secured to the facial skeleton using titanium plates and screws, typically placed at the piriform rims and zygomatic buttresses. Bone grafts may be used to fill gaps created by the repositioning.
  
  \item \textbf{Closure:} The intraoral incisions are closed with dissolvable sutures. A nasal pack may be placed temporarily, and light intermaxillary fixation may be applied to guide the bite during initial healing.
\end{itemize}

\section*{Preoperative Workup \& Preparation}
A thorough preoperative workup is essential for the success of the Le Fort I osteotomy, involving a multidisciplinary team.

\begin{itemize}
  \item \textbf{Diagnostic Records:} Comprehensive records, including panoramic and cephalometric X-rays, dental models, and photographs, are essential for assessing skeletal and dental relationships.
  
  \item \textbf{Virtual Surgical Planning:} 3D virtual planning allows for simulation of the osteotomy and design of custom surgical guides.
  
  \item \textbf{Orthodontic Preparation:} Pre-surgical orthodontics is necessary to align the teeth and prepare the dental arches for surgery.
  
  \item \textbf{Medical Clearance:} A thorough medical evaluation is performed to ensure the patient is fit for surgery, including blood tests and ECG.
  
  \item \textbf{Medication Review:} Patients are advised to stop certain medications, such as blood thinners, prior to surgery.
  
  \item \textbf{Smoking Cessation:} Patients are encouraged to stop smoking to improve healing outcomes.
  
  \item \textbf{NPO Status:} Patients must adhere to fasting instructions prior to surgery.
  
  \item \textbf{Informed Consent:} A detailed discussion of the procedure, risks, and benefits is conducted to obtain informed consent.
  
  \item \textbf{Pre-operative Antibiotics:} Prophylactic antibiotics are administered to reduce the risk of infection.
\end{itemize}

\section*{Contraindications \& Precautions}
\textbf{Absolute Contraindications:}
\begin{itemize}
  \item Uncontrolled systemic medical conditions that increase surgical risk.
  \item Active infections in the surgical field.
  \item Unrealistic patient expectations or psychological instability.
  \item Incomplete skeletal growth in adolescents.
  \item Active malignancy in the surgical area.
  \item Severe coagulopathy.
\end{itemize}

\textbf{Relative Contraindications and Precautions:}
\begin{itemize}
  \item Smoking, which delays healing and increases infection risk.
  \item Poor oral hygiene or untreated periodontal disease.
  \item Certain medications that impair bone healing.
  \item Previous radiation therapy to the head and neck.
  \item Significant temporomandibular joint dysfunction.
  \item Pregnancy, where elective surgery is deferred.
  \item Severe obstructive sleep apnea requiring careful management.
\end{itemize}

\section*{Complications \& Risks}
While the Le Fort I osteotomy is generally safe, potential complications can arise, categorized as general surgical risks and procedure-specific risks.

\textbf{General Surgical Risks:}
\begin{itemize}
  \item \textbf{Anesthesia Complications:} Risks associated with anesthesia, including respiratory issues and cardiovascular events.
  \item \textbf{Bleeding:} Intraoperative or postoperative hemorrhage may occur.
  \item \textbf{Infection:} Surgical site infections or hardware-related infections may necessitate further intervention.
  \item \textbf{Swelling and Bruising:} Significant swelling is common and typically resolves over weeks.
  \item \textbf{Pain:} Postoperative pain is expected but manageable with analgesics.
\end{itemize}

\textbf{Procedure-Specific Risks:}
\begin{itemize}
  \item \textbf{Nerve Injury:} Potential for temporary or permanent sensory changes.
  \item \textbf{Dental Injury:} Damage to adjacent teeth may occur.
  \item \textbf{Relapse:} The maxilla may shift back toward its original position.
  \item \textbf{Non-union or Malunion:} Failure of bone segments to heal correctly.
  \item \textbf{Sinus Complications:} Risks of sinusitis or oroantral fistula.
  \item \textbf{Nasal Changes:} Alterations in nasal shape may require secondary procedures.
  \item \textbf{Hardware Complications:} Issues with fixation plates or screws.
  \item \textbf{TMJ Dysfunction:} New or worsening TMJ symptoms may occur.
  \item \textbf{Speech Changes:} Temporary alterations in speech may be noted.
  \item \textbf{Vision Changes:} Rare potential for vision changes due to orbital injury.
\end{itemize}

\section*{Postoperative Recovery Timeline}
Following the Le Fort I osteotomy, patients are closely monitored in the Post-Anesthesia Care Unit (PACU) for vital signs and airway patency. Significant facial swelling and discomfort are expected, and pain management is initiated with analgesics. 

Patients typically remain hospitalized for 1-2 days, during which intravenous fluids are administered, and a liquid or soft diet is initiated. Light intermaxillary fixation may be applied to stabilize the bite. Over the first 1-2 weeks, swelling gradually subsides, and patients are restricted to a non-chewing diet for 4-6 weeks to facilitate healing. 

Physical activity is limited, and patients are advised to avoid strenuous activities for several months. Regular follow-up appointments are scheduled to monitor healing and adjust orthodontic treatment as needed. Full return to normal activities and diet occurs gradually, with complete bone healing taking up to a year.

\section*{Surgical Findings \& Pathology}
During the Le Fort I osteotomy, the surgeon documents the specific malposition of the maxilla, including any asymmetries or deficiencies. Key intraoperative findings include the integrity of the infraorbital nerves and vascular structures. While no diseased tissue is typically resected, any incidental findings, such as suspicious lesions, may be sent for histopathological examination. The primary focus is on correcting the skeletal malformation and ensuring proper stabilization of the repositioned maxilla.

\section*{Expected Postoperative Course}
A successful postoperative course following a Le Fort I osteotomy is characterized by a gradual reduction in swelling and discomfort. Pain typically resolves within 1-2 weeks, and sensory changes in the upper lip and cheeks improve over several months. Patients progress from a liquid to a soft diet over 4-6 weeks, with a gradual return to normal eating patterns as healing progresses. 

The occlusion should feel stable, and patients can resume light activities within a few weeks. Full return to strenuous activities is usually delayed until complete bone consolidation, typically 3-6 months post-surgery. The ultimate goal is to achieve a stable, functional occlusion and improved facial aesthetics.

\section*{Postoperative Care \& Follow-up}
Postoperative care is crucial for monitoring healing and ensuring the long-term success of the procedure.

\begin{itemize}
  \item \textbf{Surgical Follow-up Visits:} Initial visits are scheduled at 1 week, 2-4 weeks, 3 months, 6 months, and 1 year post-surgery.
  
  \item \textbf{Orthodontic Adjustments:} Orthodontic treatment resumes a few weeks after surgery to fine-tune the bite.
  
  \item \textbf{Imaging Studies:} Panoramic X-rays are taken at 3 and 6 months to assess bony union.
  
  \item \textbf{Elastic Wear:} Patients may be instructed to wear elastic bands to guide the bite.
  
  \item \textbf{Oral Hygiene Instruction:} Emphasis on meticulous oral hygiene to prevent infection.
  
  \item \textbf{Activity Resumption:} Gradual return to normal activities is advised, with caution regarding contact sports.
  
  \item \textbf{Warning Signs:} Patients are educated on signs of complications, such as persistent fever or significant changes in sensation.
\end{itemize}

Completion of orthodontic treatment marks the final step in the overall orthognathic treatment plan.

\section*{Epidemiology}
Le Fort I osteotomy is performed globally to address various dentofacial deformities. The incidence of skeletal malocclusions requiring surgical intervention is estimated to affect 5\% to 15\% of the population. The typical demographic for this procedure includes adolescents and young adults, generally after skeletal growth has ceased. 

Common indications include Class II malocclusion with maxillary retrusion and Class III malocclusion with maxillary deficiency. The overall morbidity associated with the procedure is low, with serious complications being rare. Minor complications, such as temporary nerve paresthesia and swelling, are more frequent but typically resolve without significant intervention.

\section*{Outcomes \& Prognosis}
The outcomes of Le Fort I osteotomy are generally favorable, with high success rates in achieving functional and aesthetic improvements. Studies indicate that over 90\% of patients report satisfaction with their facial appearance and improved oral function. 

The prognosis for long-term stability of the maxillary position is excellent, particularly with the use of rigid internal fixation and diligent post-surgical orthodontic management. Relapse rates are low, typically ranging from 5\% to 15\%, and are often manageable with orthodontic adjustments. Factors influencing prognosis include the magnitude of surgical movement, fixation stability, and patient compliance with postoperative instructions. Overall, the Le Fort I osteotomy provides a reliable solution for correcting maxillary discrepancies, leading to significant improvements in both function and aesthetics.

\section*{References}
\begin{enumerate}
  \item Hupp JR, Ellis E, Tucker MR. Contemporary Oral and Maxillofacial Surgery. 7th ed. Elsevier; 2019.
  
  \item Fonseca RJ, Marciani RD, Turvey TA. Oral and Maxillofacial Surgery. 3rd ed. Elsevier; 2017.
  
  \item American Association of Oral and Maxillofacial Surgeons. Parameters of Care: Clinical Practice Guidelines for Oral and Maxillofacial Surgery. 6th ed. AAOMS; 2021.
  
  \item Proffit WR, White RP Jr, Sarver DM. Contemporary Orthodontics. 6th ed. Elsevier; 2019.
\end{enumerate}

\clearpage